
\begin{table}[!tbp]
\centering
\caption{Load benchmarks, evidence status, and evaluation contracts.}
\label{tab:contracts}
\scriptsize
\renewcommand{\arraystretch}{0.90}
\begin{tabular}{p{0.17\textwidth}p{0.19\textwidth}p{0.13\textwidth}p{0.13\textwidth}p{0.16\textwidth}}
\toprule
Cohort & Primitive region & Evaluation units & Feedback & Role \\
\midrule
NYISO exposed development & Hourly joint region, 11 zones & 7 origins, 2018--2025 & Strict lag-2 & Primary benchmark \\
GEFCom2014 replay & Daily joint region, 24 hours & 5 quarterly origins & Strict lag-2 & Transfer sensitivity \\
CAISO development replay & Daily joint region, 24 hours & 12 origins & Strict lag-2 & Transfer development \\
CAISO post-split replay & Daily joint region, 24 hours & 6 origins & Strict lag-2 & Exposed transfer sensitivity \\
NYISO recovery & Hourly joint region, 11 zones & 181 days, 12 blocks & Strict lag-2 & Recovery-confirmatory \\
\bottomrule
\end{tabular}
\par\footnotesize Primary NYISO and recovery action paths use production shortfall-sum DRB. GEFCom/CAISO allocation sensitivities use the separate thresholded-count variant (Table~\ref{tab:transfer}); no-retune production replay is reported separately in Supplementary Table~S10. NYISO recovery retained a frozen analysis after technical recovery with prior outcome access. Cohorts and coverage geometries are not pooled.
\end{table}
